Gostaria de expressar minha sincera gratidão a todos que, de alguma maneira, contribuíram para a realização deste trabalho e para o meu crescimento acadêmico e pessoal.

Primeiramente, agradeço profundamente à minha família, que me ofereceu apoio durante toda essa jornada. Em especial, à minha mãe, que esteve sempre ao meu lado, oferecendo amor, compreensão e força nos momentos mais desafiadores.

Aos meus amigos, que tornaram este caminho mais leve e prazeroso, minha eterna gratidão. Em particular, ao meu namorado Artur, que sempre acreditou em mim, me incentivou e esteve presente, proporcionando o apoio emocional que tanto precisei. Seu carinho e confiança foram essenciais. Amo você.

Agradeço também ao meu orientador, Prof. Michael, cuja orientação técnica e apoio constante foram indispensáveis para o desenvolvimento deste trabalho. Sua competência profissional, aliada à sua amizade e à maneira generosa de compartilhar seu conhecimento, foram fundamentais para meu aprendizado e para a conclusão desta dissertação.

Sou muito grata aos meus colegas do laboratório, cujos momentos de troca e convivência tornaram a experiência acadêmica ainda mais enriquecedora. Igor, Robert, Márcio, Miguel, Artur, Hellen, Letícia, Lucas, Alessandro e Thiago, obrigado por todo o apoio, pelas conversas e pelas risadas que tornaram o ambiente de trabalho mais agradável e estimulante.

Gostaria de agradecer também ao Prof. Ascanio, que, além de disponibilizar um laboratório de excelência, me ofereceu o suporte técnico essencial para a realização deste estudo. Sua orientação foi crucial, e sou muito grata pela oportunidade de contar com seu conhecimento.

Finalmente, agradeço à FUNCAP pela concessão da bolsa de mestrado, que me possibilitou dedicar-me integralmente a este projeto.

A todos vocês, o meu mais sincero obrigado. Este trabalho é fruto de um esforço coletivo, e sem o apoio de cada um, não teria sido possível chegar até aqui. Muito obrigada!